\chapter{基于多智能体近端策略优化的有限视域多传感器协同搜索与跟踪方法}

本章的重点在于解决动态环境下传统优化方法计算复杂度高且缺乏长时序规划能力的问题，
实现多传感器协同搜索与跟踪的智能化决策。 文中首先构建了多传感器协同决策模型，
并将求解方案归纳为观测空间重构、神经网络架构设计及多智能体强化学习训练三个核心环节。
在观测层面上，本文充分利用第三章中 GM-PHD 滤波器的后验信息，通过将变长的目标高斯分量映射为固定维度的后验强度信息，
实现了观测空间从时变特征向标准化输入的转换。在决策模型上，基于集中式训练分布式执行（CTDE）框架，
设计了融合卷积神经网络（CNN）的 Actor-Critic 结构以提取高维空间特征。
针对搜索与跟踪任务的多目标性，本文进一步构建了联合奖励函数以引导策略学习。
最后通过仿真实验，验证所提方法在非线性场景下的决策优势与实时性。

\section{多智能体协同决策建模}
鉴于环境的动态性与传感器感知的局部性，我们将多传感器协同协同决策模块建模为一个部分可观测的多智能体协同决策过程。
为了方便描述，我们将该过程定义为一个六元组$\langle\mathcal{N},\mathcal{S},\mathcal{O},\mathcal{A},\mathcal{R},\gamma\rangle$
各要素定义如下：

智能体集合$\mathcal{N}$：包含$\ N$个同构的移动传感器节点。

状态空间$\mathcal{S}$:全局状态$s_{k}={s_{k}^{1},s_{k}^{2},\ldots,s_{k}^{N}}\in\mathcal{S}$
包含所有智能体的状态以及监视区域内所有目标的真实状态。

局部观测空间$\mathcal{O}$:由于传感器的视域受限，智能体无法直接获取全局状态$s_{k}$,第$i$个智能体
在$k$时刻的观测$o_{i,k}\in\mathcal{O}$包含了其自身的运动状态以及有限视域下对目标的估计状态。

联合动作空间$\mathcal{A}$:所有智能体的联合状态空间$\mathcal{A}=\mathcal{A}^{1}\times\mathcal{A}^{2}\times\cdots\times\mathcal{A}^{n}$
，其中，联合动作$\mathbf{a}=\{a^{1},a^{2},\cdots,a^{N}\}\in\mathcal{A}$。

奖励函数${\mathcal{R}}$：${\mathcal{R}}(s,\mathbf{a})\to\mathbb{R}$表示在状态$s$下执行联合动作$a$
折扣因子$\gamma$：衡量未来回报的重要性。

多传感器协同协同决策模块的最终目标是是寻找一个最优联合策略参数$\pi^\ast$，使得多智能体在时间$T$内的累积期望折扣回报最大化。记为：
\begin{equation}
J(\pi)=\mathbb{E}_{\tau\sim\pi}\left[\sum_{t=0}^T\gamma^tr_t\right]
\end{equation} 
\begin{equation}
\pi^{*}=\underset{\pi}{\operatorname*{\mathrm{argmax}}}J(\pi)
\end{equation}
其中，$\tau$表示由策略$\pi$产生的状态-动作轨迹序列，通过最大化该目标函数，系统能够在动态环境下实现对监视区域内多目标的搜索与跟踪。

\subsection{基于GM-PHD后验信息的观测空间设计}
由于传感器的视域有限，每个传感器自身视域内的观测信息，
本文将局部观测空间$O$设计为包含传感器自身状态${o}_{s}$和视域内目标状态特征的集合$O{grid}$。
鉴于随机有限集框架下目标数量的时变性，智能体直接以目标状态集作为策略网络输入存在维度不匹配的问题。
本文采用一种基于极坐标网格化的特征张量构建方法，将PHD滤波器输出的后验强度转换为固定维度的张量输入。
对强度函数进行扇形区域划分与特征提取的合理性主要源于以下两方面：
其一，在随机有限集理论中，强度函数实质上是目标状态集的一阶矩，其在任意状态空间子集上的积分，
在物理意义上严格等价于该区域内目标的期望个数。
\begin{equation}
\widehat{N}_{S}=\int_{S}v_{k}(\mathbf{x})d\mathbf{x}
\end{equation}
因此，将连续的局部观测空间离散化为若干互不重叠的扇形子区域，并聚合落入其中的高斯分量特征，
能够在保证策略网络输入维度匹配的同时有效保留目标在局部的空间分布密度。
其二，考虑到智能体搭载的主流传感器通常基于距离和方位角进行探测，
其物理视域和分辨单元天然呈现为极坐标下的扇形结构。采用极坐标网格对视域进行划分，
比直角坐标系更能精确契合传感器的实际物理探测机理与观测边界


极坐标网格化的特征张量构建构建过程如下：

(1)空间映射与网格化

假设第$i$个智能体在$k$时刻的位置为${p}_{i,k}=\begin{bmatrix}x_{i,k},y_{i,k}\end{bmatrix}^{T}$，以${p}_{i,k}$为原点，
航向角$\theta_{i,k}$为极轴正方向，建立局部坐标系。将视域沿着极径和极角的方向进行离散化，
划分为$U\times{V}$个网格单元，其中$U$和$V$分别为沿着极径和极角方向的网格划分数量。
每个网格记为$g(\mathrm{u,v})$，其中$u\in[0,U-1]$表示极径索引，$v\in[0,V-1]$表示极角索引。
设第$j$个目标在全局坐标系下的位置坐标为$(m_{x,j},m_{y,j})$，将其转换至局部极坐标系的计算过程如下。
目标相对于智能体的距离$\rho_j$为：
\begin{equation}
\rho_j=\left\| \mathbf{m}_j - \mathbf{p}_{i,k} \right\|_2
\end{equation}
目标相对于智能体的极角$\theta_j$为：
\begin{equation}
\theta_j=\arctan\left(m_{y,j}-p_{y,i,k},m_{x,j}-p_{x,i,k}\right)-\theta_{i,k}
\end{equation}
其中，$\theta_{i,k}$表示智能体$i$在$k$时刻的航向角；
在此映射下，相对极角$\theta_j$的值域被规范化至$[-\pi,\pi)$区间。

(2)特征张量构建

在完成高斯分量到局部极坐标系的空间映射后，需要进一步将高斯粒子集合转化为策略网络可直接接收的结构化数据。
该构建过程的核心在于对网格内的粒子属性（如权重、均值、协方差等）进行多维度特征提取与通道聚合。
基于上述网格映射，本文构建特征张量$\mathbf{O}\in\mathbb{R}^{\mathbf{U}\times\mathbf{V}\times\mathbf{C}}$,
$c$为通道数，在本文构建的极坐标网格中，每个网格单元$g(\mathrm{u,v})$对应物理空间中的一个特定扇环区域，
定义集合$s_{u,v}$为落入该网格单元$g(\mathrm{u,v})$的所有高斯分量的索引集合。各通道的特征提取方式如下：

强度通道：特征张量在概率权重通道的取值定义为网格内所有高斯分量权重之和：
\begin{equation}
\mathbf{0}(u,v,1)=\sum_{j\in\mathcal{S}_{u,v}}w_{j}
\end{equation}

速度特征表示目标接近或远离传感器的趋势。设网格单元$g(\mathrm{u,v})$内第$j$个高斯分量的均值向量
$\mathbf{m}_{j}=\left[m_{x,j},m_{vx,j},m_{y,j},m_{vy,j}\right]^{T}$,设$\mathbf{v}_{j}=(m_{vx,j},m_{vy,j})$，
智能体航向角为$\theta_s$，定义由全局坐标系旋转至航向坐标系的旋转矩阵为：
\begin{equation}
R(\theta_s)=\begin{bmatrix}\cos\theta_s&\sin\theta_s\\-\sin\theta_s&\cos\theta_s\end{bmatrix}
\end{equation}
则该高斯分量在航向对齐坐标系下的速度表示为：
\begin{equation}
\tilde{\mathbf{v}}_j=\begin{bmatrix}\widetilde{m}_{vx,j}\\\widetilde{m}_{vy,j}\end{bmatrix}=R(\theta_s)\mathbf{v}_j
\end{equation}
其中$\tilde{m}_{vx,j}$表示沿智能体航向方向的速度分量，$\tilde{m}_{vy,j}$表示垂直于航向方向的横向速度分量。在每个网格单元
$g(\mathrm{u,v})$内，对所有落入该区域的高斯分量速度进行权重归一化平均，则网格单元$g(\mathrm{u,v})$的速度特征分别为：
\begin{equation}
O(u,v,2)=\frac{\sum_{j\in S_{u,v}}w_{j}\widetilde{m}_{vx,j}}{\sum_{j\in S_{u,v}}w_{j}+\epsilon}
\end{equation}
\begin{equation}
O(u,v,3)=\frac{\sum_{j\in\mathcal{S}_{u,v}}w_{j}\widetilde{m}_{vy,j}}{\sum_{j\in\mathcal{S}_{u,v}}w_{j}+\epsilon}
\end{equation}
其中$\epsilon$为防止分母为零的极小常数。
目标状态协方差的对角线反映了目标状态的分散程度，因此将网格单元$g(\mathrm{u,v})$内所有高斯分量
位置对角线之和表征其不确定性特征。
\begin{equation}
\mathbf{O}(u,v,4)=\begin{cases}\frac{\sum_{j\in\mathcal{S}_{u,v}}w_j\cdot\mathrm{tr}\left(W\mathbf{P}_jW^T\right)}{\sum_{j\in\mathcal{S}_{u,v}}w_j+\epsilon},&\quad\mathcal{S}_{u,v}\neq\emptyset\\0,&\quad\mathcal{S}_{u,v}=\emptyset\end{cases}
\end{equation}
其中，$W$是用来提取协方差矩阵中有关目标位置维分量的合适矩阵。

动作空间

为简化控制模型，假设智能体动作空间是离散且同构的，智能体的控制输出定义为航向角的调整量。
第$\mathrm{i}$个传感器节点的离散动作的角度值$a_{i}$可以表示为：
\begin{equation}
a_i=-\pi+\frac{2\mathrm{n_a}\pi}{\mathrm{N_a}},\quad\mathrm{n_a\in[0,N_a-1]}
\end{equation}
其中，$N_a$表示动作空间离散化后的动作集合，$n_a$表示离散动作的动作索引。

\subsection{搜索与跟踪的复合奖励函数设计}
本文中奖励函数包含跟踪精度奖励、新目标发现奖励及协同与约束奖励三部分，旨在优先保障高价值目标连续跟踪的前提下，
最大化系统对未知环境的探索能力。
跟踪精度效能奖励：

对于已经检测到的目标，希望尽量减少目标状态估计的不确定性，为了鼓励智能体提升对已发现目标的估计精度。
因此考虑采用所有已检测到的目标位置协方差的对角线总和作为跟踪性能奖励函数。
但是由于智能体为了不受到跟踪精度的惩罚，就放弃对已发现目标的跟踪，所以加入与目标基数估计
\begin{equation}
r_{track}=-\sum_{i\in\mathbb{I}_k}\mathrm{tr}\left(WP_k^iW^T\right)+\beta(N_{k+1}-N_k)
\end{equation}
$\mathbb{I}_k$表示$k$时刻已被检测并正在跟踪的目标集合,$P_k^i$表示第$i$个目标在时刻$k$的状态协方差，
$W$表示提取位置维度的选择矩阵。$N_{k+1}$和$N_k$分别表示$k+1$时刻和$k$时刻的目标基数估计。

新发现目标奖励：

当系统当前时刻检测到的确认目标总数超过过去一段时间内的最大目标数时，视为发现了新目标，
给予固定的正向奖励 $R_{new}$，以激励智能体探索未观测区域。
\begin{equation}
r_{new} = 
\begin{cases}
\alpha N_k, & N_k > \max \{N_{k-L}, N_{k-L+1}, \dots, N_{k-1}\} \\
0, & N_k \leq \max \{N_{k-L}, N_{k-L+1}, \dots, N_{k-1}\}
\end{cases}
\end{equation}
其中，$\alpha$为奖励因子，$N_k$表示$k$时刻的目标基数估计，$L$表示滑动的历史时间窗口。

视域重叠度奖励：

防止传感器网络出现覆盖盲区或因过度聚集导致的资源浪费，引入视域重叠度奖励$r_{\mathrm{ol}}^{i}$
该奖励机制旨在引导智能体维持一个理想的视域重叠比例。设$F_i$、$F_j$分别为智能体$i$、$j$的视域覆盖区域，
智能体$i$的视域重叠率$\rho_i$表示为：
\begin{equation}
\rho_i=\frac{\left|F_i\cap\left(\bigcup F_{\mathrm{j}}\right)\right|}{\left|F_i\right|}
\end{equation}
智能体i的视域重叠度奖励函数表示为：
\begin{equation}
r_{\mathrm{ol}}^i=\delta \cdot \rho_i
\end{equation}
其中，$\delta$为奖励权重。

边界惩罚：

传感器太靠近边界会造成部分观测区域会落在监测区域外，造成观测资源浪费，边界惩罚可描述如下：
\begin{equation}
r_{bound}=\begin{cases}-0.5\cdot(d-d_{i}^{b})/d,&\quad d_{i}^{b}\leq d \\ 0,&\quad d_{i}^{b}>d\end{cases}
\end{equation}
其中，$d_{i}^{b}$为智能体$i$距离监测区域边界的最短距离。$d$为传感器的视域半径。

综上所述，智能体的总奖励函数定义为：
\begin{equation}
R=\lambda_1r_{track}+\lambda_2r_{new}+\lambda_3r_{ol}+\lambda_4r_{bound}
\end{equation}
其中，$\lambda_1,\lambda_2,\lambda_3,\lambda_4$均为奖励因子权重。

\section{基于多智能体近端策略优化的协同搜索与跟踪控制算法}
根据前文建立的多传感器协同决策与离散动作空间模型，可以采用深度强化学习方法来学习智能体
的协同搜索与跟踪动作策略。在众多多智能体强化学习算法中，多智能体近端策略优化
（Multi-Agent Proximal Policy Optimization, MAPPO）算法作为一种先进的 Actor-Critic 架构算法，
在处理复杂的多目标博弈与协同管控问题时展现出了卓越的性能。相比于其他策略梯度算法，
MAPPO 算法具有超参数鲁棒性强、训练稳定性高且不易陷入局部最优等显著优点。因此，本节结合协同管控模型与
特征张量表示方法，设计了用于学习智能体集群协同搜索与跟踪策略的MAPPO算法。

在多智能体系统进行协同决策时，每一个智能体策略网络的输入变量包括对目标的局部观测信息和局部环境信息
（如自身位置、速度和任务区域边界等）；其输出变量为动作空间上的最优策略概率分布以及当前状态的价值评估。
由于环境中出现的动态目标观测信息通常是变长的，难以直接被神经网络处理，因此需要使用特征张量表示方法，
将这些非结构化、变长的数据统一映射为固定维度的特征张量，随后将特征张量输入到神经网络中进行计算。
据此，设计的 MAPPO 算法网络结构如图 4.1 所示。其中，特征张量表示结果使用卷积神经网络（CNN）
进一步提取空间高维特征，其它低维的自身状态信息使用全连接层（FC）进行处理；
处理后的特征向量经串联操作后，分别输入至 Actor 网络输出动作策略，
以及输入至 Critic 网络评估状态值函数，从而指导整个智能体集群完成高效的协同搜索与跟踪任务。

在网络参数更新阶段，MAPPO 算法通过广义优势估计（GAE）计算优势函数$\hat{A}_t$,以衡量当前动作相较于平均策略的优劣。
为了保证策略训练的平稳性，防止单次更新步长过大导致策略崩溃，Actor 网络参数$\theta$的目标函数$L(\theta)$引入了裁剪函数
\begin{equation}
L(\theta) = \mathbb{E}{t}\left[ \min \left( \rho_t(\theta)\hat{A}t, \text{clip}\left( \rho_t(\theta), 1-\epsilon, 1+\epsilon \right)\hat{A}t \right) \right]
\end{equation}
其中，$\rho_t(\theta) = \frac{\pi\theta(a_t|o_t)}{\pi{\theta{old}}(a_t|o_t)}$ 表示新旧策略的比率；
$\epsilon$ 为截断超参数，用于限制策略更新幅度。

Critic 网络参数$\phi$的优化目标则是最小化全局状态价值预测值与实际多步折扣回报$\hat{R}_t$之间的均方误差，
其损失函数$L(\phi)$定义为：
\begin{equation}
L(\phi) = \mathbb{E}{t}\left[ \left( V\phi(s_t) - \hat{R}_t \right)^2 \right]
\end{equation}

\section{算法整体流程}
综上所述，本文基于特征张量与 MAPPO 的有限视域多传感器协同算法的完整执行流程可归纳为以下步骤：

\begin{table}[htbp]
  \centering
  %\caption{基于特征张量与 MAPPO 的有限视域多传感器协同管控算法}
  \label{alg:mappo_coop}
  \renewcommand{\arraystretch}{1.2} % 稍微增加行高使得公式显示更美观
  \begin{tabular}{p{0.95\textwidth}}
    \toprule
    \textbf{Algorithm 4.1} 基于特征张量与 MAPPO 的有限视域多传感器协同管控算法 \\
    \midrule
    \textbf{初始化}: Actor 策略网络参数 $\theta$ 和全局 Critic 价值网络参数 $\phi$；经验缓存 $\mathcal{D}$ \\
    
    \textbf{// 集中式训练:} \\
    1: \textbf{for} $epi = 1$ \textbf{to} $episodes$ \textbf{do} \\
    2: \quad 环境与传感器状态初始化; \\

    3: \quad \textbf{for} $t = 1$ \textbf{to} $T$ \textbf{do} \\
    4: \quad \quad 各传感器进行本地滤波，经多传感器信息融合后，获得融合后的目标后验估计; \\
    5: \quad \quad 对每一个传感器 $i$，结合自身位姿与融合后的目标后验估计构建特征张量 $\mathbf{O}_{i,t}$，并基于 $\pi_{\theta}(\cdot|\mathbf{O}_{i,t})$ 采样动作 $a_{i,t}$; \\
    6: \quad \quad 执行联合动作 $\mathbf{a}_t = \{a_{1,t}, \cdots, a_{N,t}\}$，观测环境获取新量测，计算即时奖励 $r_t$; \\
    7: \quad \quad 将全局经验 $(s_t, \mathbf{O}_t, \mathbf{a}_t, r_t, s_{t+1})$ 存储在经验缓存 $\mathcal{D}$ 中; \\
    8: \quad \textbf{end for} \\
    
    9: \quad 从经验缓存 $\mathcal{D}$ 中提取本回合完整的轨迹序列数据; \\
    10: \quad 利用 Critic 网络预测的全局价值 $V_\phi(s_t)$ 和广义优势估计 (GAE)，计算单步优势函数 $\hat{A}_t$ 及目标回报 $\hat{R}_t$; \\
    11: \quad 执行梯度上升算法最大化截断代理目标函数，更新所有智能体共享的 Actor 网络参数 $\theta$ \\
    12: \quad 执行梯度下降算法最小化均方误差损失函数，更新全局 Critic 网络参数 $\phi$ \\
   
    13: \quad 清空经验缓存 $\mathcal{D}$; \\
    14: \textbf{end for} \\
    
    \textbf{// 分布式执行:} \\
    15: 环境与传感器状态初始化; \\
    16: 为所有传感器节点加载 Actor 策略网络; \\
    17: \textbf{for} $t = 1$ \textbf{to} $T$ \textbf{do} \\
    18: \quad 各传感器进行本地滤波，经多传感器信息融合后，获得融合后的目标后验估计; \\
    19: \quad 对每一个传感器 $i$，结合自身位姿与融合后的目标后验估计构建局部特征张量 $\mathbf{O}_{i,t}$; \\
    20: \quad 智能体 $i$ 利用本地 Actor 网络输出确定性最优动作 $a_{i,t} = \arg\max_a \pi_{\theta^*}(a|\mathbf{O}_{i,t})$; \\
    21: \quad 执行联合动作 $\mathbf{a}_t$，更新环境状态; \\
    22: \textbf{end for} \\
    \bottomrule
  \end{tabular}
\end{table}



\section{仿真实验与结果分析}
\section{本章小结}

